\documentclass[12pt,a4paper]{article}
\usepackage[left=7.6667mm,right=7.3333mm,top=6.3333mm,bottom=6.3333mm,footskip=3mm]{geometry}
\usepackage{fontspec,amsmath,booktabs}
\usepackage[russian]{babel}
\setmainfont{Times New Roman}
\usepackage[hidelinks]{hyperref}
\hypersetup{pdftitle={Убыточные периоды шести стратегий на BTC}}
\setlength{\parindent}{0pt}
\setlength{\parskip}{3pt}
\setlength{\emergencystretch}{2em}
\newcommand{\strategy}[1]{\par\addvspace{7pt}\noindent{\large\bfseries #1}\par\nobreak}
\begin{document}
\pagestyle{plain}
{\LARGE\bfseries Убыточные периоды шести стратегий на BTC}\par
Deribit, 1 января 2024~г. --- 14 ноября 2026~г. Денежный результат учитывает модельные издержки. Просадки рассчитаны по дневным значениям капитала.

\strategy{1. Чистый стрэддл}
Весь 2024 год счёт оставался в деньгах: доходность 0\%, первая покупка --- 22 апреля 2025~г. Убыток возник в двух августовских циклах. Для $Q$ пар со страйком $K$, премией $P$ на пару и расходами $C$ результат равен $\Pi=Q(|D_T-K|-P)-C$, где $D_T$ --- цена BTC при экспирации. Обе позиции имели объём $Q=0{,}1$. В обоих случаях выплата оказалась меньше премии ещё до расходов.

\begin{center}\begin{tabular}{@{}lrrr@{}}\toprule
Цикл, 2025~г. & Премия пары, USDC & Выплата пары, USDC & Результат позиции, USDC\\\midrule
12--13 августа & 2\,068,23 & 1\,383,17 & $-77{,}39$\\
15--16 августа & 763,33 & 712,02 & $-12{,}56$\\\bottomrule
\end{tabular}\end{center}

\strategy{2. Стрэддл + грид}
В 2024~г. входов не было. На тех же августовских входах $\Pi_{\text{связки}}=\Pi_{\text{стрэддла}}+\Pi_{\text{грида}}$. 12--13 августа сетка добавила убыток: $-77{,}39-0{,}87=-78{,}26$ USDC. 15--16 августа вклад сетки равнялся нулю, итог составил $-12{,}56$ USDC.

В исходном эксперименте задания~3 (май 2022~г. --- ноябрь 2026~г.) тот же эффект виден в 19 боковых циклах. Стрэддл потерял 187\,886 USDC, сетка заработала 3\,466 USDC. Торговля колебаниями возместила лишь $3\,466/187\,886\approx1{,}8\%$ опционного убытка.

\strategy{3. Самостоятельный грид}
Доходность составила $-0{,}12\%$ в 2024~г. и $-0{,}06\%$ в 2025~г. Худший цикл --- 3 февраля 2025~г. Бот продал 0,08 BTC по 92\,236 USDC и выкупил по 93\,993. Результат шорта до расходов равнялся $0{,}08(92\,236-93\,993)=-140{,}56$ USDC, чистый убыток --- $142{,}80$ USDC. Цена обратной покупки превысила цену продажи.

\strategy{4. Стрэддл + частичный дельта-хедж}
В 2024~г. входов не было. Бот держит позицию перпетуала около $H^*=-0{,}5Q\Delta$, где $\Delta$ --- суммарная дельта одной пары. 12--13 августа 2025~г. хедж заработал $8{,}38$ USDC при опционном убытке $77{,}39$ USDC: итог $-69{,}01$ USDC. 15--16 августа обе части потеряли деньги: $-12{,}56-7{,}95=-20{,}51$ USDC. Частичный хедж уменьшает чувствительность к цене BTC, но возврат премии зависит от выплаты и результата сделок хеджа.

\strategy{5. Тренд за 42 дня}
Крупнейшая дневная просадка от пика --- $31{,}19\%$ с 14 марта по 11 октября 2024~г. Длинная серия 25 мая --- 25 июня потеряла 15\,389 USDC, короткая серия 27 июня --- 27 июля --- ещё 12\,310 USDC. Весь 2024~г. завершился в плюсе; результат 2025~г. составил $-6{,}75\%$.

Знак позиции $H_t$ задаёт прошлая доходность $\ln(S_t/S_{t-42})$. Следующий ценовой результат равен $H_t(S_{t+1}-S_t)$. Когда цена разворачивается раньше сигнала, это произведение отрицательно. Последовательные развороты дают убыток в обе стороны.

\strategy{6. Тренд + стрэддл с гридом, 50/50}
С 14 марта по 11 октября 2024~г. капитал упал со 183\,824 до 149\,886 USDC, просадка --- $18{,}46\%$. Опционная часть оставалась в деньгах: $\Delta E_{\text{портфеля}}=\Delta E_{\text{тренда}}+0$. К мартовскому пику доля тренда выросла примерно до $59{,}2\%$: счёт стартовал с равными долями, затем каждая часть торговала своим капиталом. Поэтому просадка портфеля составила $0{,}592\times31{,}19\%\approx18{,}46\%$. В 2025~г. опционная часть заработала, но не покрыла потери тренда. Годовой результат портфеля --- $-2{,}48\%$.
\end{document}
